\scanpage{44}
\mspage{40}
% archive sheet number: 43

Je vais maintenant relier la relation $\unsure{\overset{\circ}{\ll}|\overset{\circ}{\Subset}}$ \inserted{(ou $\ll$)}, qui n'est pas encore intervenue dans ces considérations de supports (si ce n'est pour dire que les $S\in\Sigma_{\mathcal A}$ sont des parties $\overset{\circ}{\ll}$-fermées de $\mathcal A$), à ces notions de support. Tout d'abord, pour $\forall X\in\mathcal A$, considérons \deleted{$\unsure{X\leftarrow}$} l'ens
\[
\unsure{\mathrm{Dep}|\mathrm{D\acute{e}p}}(X)=\{\deleted{\unsure{X}}X^{\prime\circ}\mid X'\in\widetilde X\}\subset\unsure{\Sigma\mathrm{cons}|\Sigma_{\mathrm{cons}}}\subset\Sigma
\]
(``\uline{\unsure{déployement}}\deleted{'' de $X$} (``spatial'') de $X$).

On pourrait le définir pour $\unsure{H\subset|\forall}F\in\unsure{\mathscr F|\mathcal F|\mathbb F}(\mathcal A)$, mais pour l'instant on ne s'y intéresse que pour les $X\in\mathcal A$. On a

$\mathrm{Dep}(X)$ admet un plus grand élément, $X^{\circ}$, et est formé d'éléments \inserted{$\ne\varnothing_\Sigma$,} \unsure{mutuellement} disjoints ($\mathrel{|\circ|}$) de $\Sigma\mathrm{cons}$. De plus, l'application $\widetilde X\mathbin{-}\mathrm{Dep}(X)$ est bijective, et $\mathrm{Dep}(X)$ hérite de $\widetilde X$ une structure d'ordre. On considérera tjrs $\mathrm{Dep}(X)$ comme une partie de $\Sigma$, munie de \deleted{une} \inserted{sa} relation d'ordre, \unsure{cela aussi}, comme une figure (élémentaire) ensembliste dans $\Sigma$, ``\uline{séparée}'' i.e. le préordre sur son support est séparé.
